\paragraph{\trYonedaEmbedding。}\label{sub:cunital-yoneda-embedding}
按定义，对 \trCunital{}的 $\scr{A}$ 而言，\trYonedaEmbedding{}%
$\scr{J}=\scr{J}_{\scr{A}}$ 的\trImage{}落在子\trCategory{}%
$\mathit{mod}(\scr{A})\subset\mathit{nu}\text{-}\mathit{mod}(\scr{A})$ 之中。

\vkthmsync{12}
\begin{lemma}\label{thm:yoneda-evaluation-quasiisomorphism}
对任意 \trCunital{} \trAinfty-模 $\scr{M}$ 与 $\scr{A}$ 的任意\trObject{} $Y$，
\eqref{eq:yoneda-evaluation-chain-map}中的映射 $\lambda$ 是一个\trQuasiisomorphism。
\end{lemma}
\begin{proof}
按定义，$\lambda$ 的\trMappingCone{}（\trSuspension{} $-1$）是\trChainComplex
\begin{equation}\label{eq:mapping-cone-of-lambda}
\Bigl(\scr{M}(Y)\oplus\hom_{\scr{Q}}(\scr{Y},\scr{M})[-1],\;
\begin{pmatrix}
\mu^1_{\scr{M}} & 0\\
\lambda & -\mu^1_{\scr{Q}}
\end{pmatrix}\Bigr).
\end{equation}
过滤的办法是：先取子\trComplex{} $\hom_{\scr{Q}}(\scr{Y},\scr{M})[-1]$，然后再施以其通常的
长度\trFiltration。令 $A=H(\scr{A})$、$M=H(\scr{M})$，则所得的\trSpectralSequence{}具有
\begin{equation*}
E_1^{rs}=
\begin{cases}
M^s(Y), & r=0,\\[2.5ex]
\begin{aligned}[t]\displaystyle
\prod_{X_0,\ldots,X_{r-1}}\Hom^{s}_{\KK}\bigl(
&\Hom_A(X_{r-1},Y)\otimes\Hom_A(X_{r-2},X_{r-1})\otimes\cdots\\
&\cdots\otimes\Hom_A(X_0,X_1),\,M(X_0)\bigr),
\end{aligned}
& r>0.
\end{cases}
\end{equation*}
由\eqref{eq:yoneda-evaluation-chain-map}与\eqref{eq:module-composition-maps}可以
看出，\trDifferential{} $\partial=\partial^{rs}\colon E_1^{rs}\to E_1^{r+1,s}$ 在 $r=0$ 时由
$(\partial c)(b)=(-1)^{|b|}cb$ 给出，而在 $r>1$ 时则为
\begin{equation*}
\begin{split}
(\partial t)(b,a_r,\ldots,a_1)
&=(-1)^{\clubsuit_{r+|b|}}\,t(b,a_r,\ldots,a_2)a_1\\
&\quad+(-1)^{|b|}\,t(ba_r,a_{r-1},\ldots,a_1)\\
&\quad+\sum_n(-1)^{\Delta}\,t(b,a_r,\ldots,a_{n+2}a_{n+1},\ldots,a_1)
\end{split}
\end{equation*}
其中 $\Delta=|a_{n+2}|+\cdots+|a_r|+|b|-r+n+1$。这是一个标准的 \trBarResolution；当 $A$ 中的%
\trIdentityMorphism{}在 $M$ 上充当恒等时，它容有一个\trContractingHomotopy{}%
$(\kappa t)(b,a_{r-1},\ldots,a_1)=t(e_Y,b,a_{r-1},\ldots,a_1)$。于是
\eqref{eq:mapping-cone-of-lambda}是\trAcyclic{}的，这就蕴涵了所要的结论。
\end{proof}

特别地，对那些来自\trYonedaEmbedding{}的 \trAinfty-模 $\scr{M}$，我们有
$\lambda=\scr{J}^1$，于是：

\begin{corollary}\label{thm:yoneda-cohomologically-full-faithful}
$\scr{J}$ 是\trCohomology\trFullyFaithful{}的（从而也是\trCunital{}的）。
\noqed
\end{corollary}

观察 $\scr{J}$ 的\trImage，还可以看到：

\begin{corollary}\label{thm:cunital-quasiisomorphic-strictly-unital}
任何 \trCunital{} \trACategory{}都以\trCanonical{}的方式与一个\trStrictlyUnital{} dg
\trCategory{}\trQuasiisomorphism{}。
\noqed
\end{corollary}

现在我们来考虑\trYonedaEmbedding{}关于 \trCunital{} \trACategory{}之间的%
\trCohomology\trFullyFaithful{} \trAFunctor{} $\scr{F}\colon\scr{A}\to\scr{B}$
的自然性。对任意这样的\trFunctor{} $\scr{F}$，
\eqref{eq:yoneda-pullback-transformation}中那个\trNaturalTransformation{}的一阶项
是 \trCunital{} \trAinfty-模之间的一个\trIsomorphism（由%
\cref{thm:module-composition-quasiisomorphism}）。鉴于%
\cref{thm:right-composition-iso}，这蕴涵 $T$ 本身在 \trAFunctor{}的底层%
\trCohomology{}\trCategory{}中是一个\trIsomorphism。换言之，图
\begin{equation}\label{eq:yoneda-naturality-diagram}
\begin{tikzcd}
\scr{A}\ar[r,"\scr{F}"]\ar[d,"\scr{J}_{\scr{A}}"] & \scr{B}\ar[d,"\scr{J}_{\scr{B}}"]\\
\mathit{mod}(\scr{A}) & \mathit{mod}(\scr{B})\ar[l,"\scr{F}^*"]
\end{tikzcd}
\end{equation}
在 $H^0(\mathit{fun}(\scr{A},\mathit{mod}(\scr{A})))$ 中交换至相差一个%
\trCanonical{}\trIsomorphism{}的程度。

\begin{remark}\label{thm:yoneda-ext-interpretation}
下面是\cref{thm:yoneda-evaluation-quasiisomorphism}之证明的一个更抽象的重述，
其精神与\cref{thm:morita-bimod-route}一致。令 $\mathit{Mod}^0(A)$ 为从 $A$ 到\trGraded\trVectorSpace{}%
的\trContravariant{}（\trUnital{}）\trFunctor{}所成的阿贝尔\trCategory。设
$\scr{M}_0$、$\scr{M}_1$ 是 \trCunital{} \trAinfty-模，则其底层的\trCohomology{}模
$M_0$、$M_1$ 是 $\mathit{Mod}^0(A)$ 的\trObject。经由 \trBarResolution，来自
$\hom_{\scr{Q}}(\scr{M}_0,\scr{M}_1)$ 的长度\trFiltration{}的\trSpectralSequence{}，其 $E_2$
项可以解释为
\begin{equation}\label{eq:ext-e2-term}
E_2=\operatorname{Ext}_{\mathit{Mod}^0(A)}(M_0,M_1),
\end{equation}
其中右边按与\eqref{eq:hh-as-ext} 中相同的方式双\trGraded。与\trYoneda{} \trAinfty-模相伴的诸%
\trCohomology{}模是 $\mathit{Mod}^0(A)$ 的投射\trObject。因此，若 $M_0=\scr{Y}$，
则\eqref{eq:ext-e2-term}中仅有一个非零的列，而该列可以等同于 $M_1(Y)$。于是，
$\scr{M}_1(Y)$ 与 $\hom_{\scr{Q}}(\scr{Y},\scr{M}_1)$ \trQuasiisomorphism。
\end{remark}
